% !TeX root = ../main_skripsi.tex
\chapter{METODE PENELITIAN}

\section{Jenis atau Desain Penelitian}

[Jelaskan pendekatan, jenis desain, dan alasan pemilihannya berdasarkan rumusan
masalah. Uraian harus cukup rinci untuk direplikasi.]

\section{Tempat dan Waktu Penelitian}

[Bagian ini diperlukan untuk survei atau penelitian berbasis lokasi. Jelaskan
lokasi, konteks, alasan pemilihan, periode, dan durasi. Hapus bila tidak relevan.]

\section{Populasi dan Sampel Penelitian}

[Jelaskan karakteristik populasi, unit analisis, ukuran sampel, dasar penentuan
ukuran, teknik sampling, serta kriteria inklusi dan eksklusi. Jika seluruh
sasaran dijangkau, gunakan istilah sumber data atau subjek penelitian.]

\section{Definisi Operasional Variabel}

[Turunkan definisi konseptual menjadi definisi operasional, indikator, skala,
sumber data, dan cara pengukuran yang dapat diaudit.]

\section{Teknik dan Instrumen Pengumpulan Data}

[Jelaskan teknik pengumpulan, proses penyusunan instrumen, kisi-kisi, prosedur,
serta kontrol kualitas data.]

\section{Validitas dan Reliabilitas Instrumen}

[Untuk survei, jelaskan jenis validitas, metode pengujian, kriteria keputusan,
dan uji reliabilitas. Untuk tes kognitif, tambahkan indeks kesulitan, daya beda,
dan efektivitas pengecoh jika relevan. Hapus bagian jika tidak digunakan.]

\section{Teknik Analisis Data}

[Jelaskan tahapan pengolahan, model statistik atau nonstatistik, uji prasyarat,
kriteria keputusan, perangkat lunak dan versinya, serta cara menjawab setiap
rumusan masalah.]
